\section{Face recovery}

Faces to the mean parameter represent the hard constraints in a model.
Recovering the right face therefore corresponds with a correct identification of the hard constraints.

\begin{lemma}
    When generating data by a distribution $\gendistributionwith$ with mean $\genmean$ we always have
    \begin{align*}
        \datameanwith \in \faceto{\genmean} \, .
    \end{align*}
\end{lemma}
\begin{proof}
    We always have $\nonzeroof{\gendistributionwith}\prec \nonzeroof{\empdistributionwith}$ when drawing data from $\gendistributionwith$.
    Since the distribution $\gendistributionwith$ has its mean on the face $\faceto{\genmean}$, the image of its support is included in $\genfaceimset$. %% CHECK NOTATION!
    Also the image of the support of $\empdistributionwith$ is included therein and thus $\datameanwith \in \faceto{\genmean}$. %% ADD REFERENCE TO THAT LEMMA
\end{proof}

Estimating the maximum entropy distribution with the data mean, we always estimate the face parameter as $\faceto{\datamean}$.
Since $\datameanwith \in \faceto{\genmean}$ we have $\faceto{\datamean}\subset\faceto{\genmean}$.
The constraints are therefore never underestimated.



\begin{example}[Face recovery of bernoulli]

    The mean parameter is parameter of Bernoulli distribution, the distribution of $\datamean$ is the binomial distribution
    \begin{align*}
        \datamean \sim \frac{1}{\datdim} B(\datdim,\genmean)  \, .
    \end{align*}

    Face recovery:
    \begin{itemize}
        \item When $\genmean\in\{0,1\}$, then we always have face recovery and perfect recovery.
        \item When $\genmean\in(0,1)$, then we have not always face recovery. This case is investigated in the following.
    \end{itemize}

    Expectation of squared norm of the error: $\frac{1}{\datdim} \genmean \cdot (1-\genmean)$
        %% Chebyshev bound
    We have for any $\delta>0$ the chebyshev bound
    \begin{align*}
        \probat{\absof{\datamean-\genmean}>\delta}
        = \probat{\absof{\datamean-\genmean}>\delta^2}
        \leq \frac{\genmean \cdot (1-\genmean)}{\datdim \delta^2} \, .
    \end{align*}
    Thus, choosing $\delta=\genmean$, we have for any $\eta$ with probability $1-\eta$ face recovery provided that
    \begin{align*}
        \datdim \geq
        \frac{1}{\eta} \cdot \frac{1-\genmean}{\genmean} \, .
    \end{align*}

    %% Exact probability of face recovery
    To show that the bound is tight we compute the exact probability of face recovery.
    The face recovery event is the complement of all datapoints being $0$ or $1$.
    The probability is therefore
    \begin{align*}
        1- (1-\genmean)^{\datdim} - (\genmean)^{\datdim} \, .
    \end{align*}
    If $\genmean$ is very small, then this is approximately $\datdim\cdot \genmean + O({\genmean}^{2})$.
    Therefore we need $\datdim\sim\frac{1}{\genmean}$ samples for face recovery.



\end{example}